% Ch17 WS1 -- entropy, its sign, and the second law. Blocks 1-2.
\documentclass{apchem-worksheet}

\apcunitword{Chapter}
\apcunit{17}
\apcunittitle{Entropy and Free Energy}
\apcdoctitle{Worksheet 1 \textbullet\ Entropy and the Second Law}
\apcsubtitle{Zumdahl \S17.1--17.2 \textbullet\ count moles of gas before you do anything else}

\begin{document}
\wsheader[Write \textbf{thermodynamically favored}, not ``spontaneous.''
$S^\circ$ in \si{\joule\per\kelvin\per\mole}: \ce{H2} 131 \textbullet\
\ce{O2} 205 \textbullet\ \ce{N2} 192 \textbullet\ \ce{H2O(l)} 70
\textbullet\ \ce{H2O(g)} 189 \textbullet\ \ce{CO2} 214 \textbullet\
\ce{CO} 198 \textbullet\ \ce{C(gr)} 6 \textbullet\ \ce{NH3} 193.]

\problem[]{Entropy measures dispersal. Name the two things that disperse,
and give an example of each.
\answerlines[3]{\textbf{Matter} --- particles occupy a larger volume or move
more freely, as when a solid melts or a gas expands. \textbf{Energy} --- the
distribution of molecular kinetic energies broadens as temperature rises, so
the energy is spread over more available states.}}

\problem[]{Give the sign of $\Delta S^\circ$ for each change and a
one-phrase reason:}
\begin{parts}
  \item \ce{H2O(s) -> H2O(l)}
        \answerblank[6.4cm]{positive --- solid to liquid}
  \item \ce{2SO2(g) + O2(g) -> 2SO3(g)}
        \answerblank[6.4cm]{negative --- 3 mol gas to 2}
  \item \ce{NH4NO3(s) -> NH4+(aq) + NO3^-(aq)}
        \answerblank[8.4cm]{positive --- ordered lattice to free ions}
  \item A gas is compressed into half its volume
        \answerblank[7cm]{negative --- less space to occupy}
  \item \ce{CaCO3(s) -> CaO(s) + CO2(g)}
        \answerblank[7.8cm]{positive --- gas produced from a solid}
  \item A sample of argon is warmed from 300 to \SI{400}{\kelvin}
        \answerblank[8.8cm]{positive --- energy spread over more states}
  \item \ce{H2(g) + Cl2(g) -> 2HCl(g)}
        \answerlines[2]{Near zero --- 2 mol of gas becomes 2 mol of gas, so
        $\Delta n_{\text{gas}} = 0$ and no large change is expected. (It is
        slightly positive, $+20$~J/K, because \ce{HCl} has more atoms free
        to arrange than the symmetric diatomics.)}
\end{parts}

\problem[]{Explain why counting moles of gas is the first move, using
water as evidence. \answerlines[3]{Gases carry far more entropy than
condensed phases, so a change in the number of moles of gas dominates any
other contribution. \ce{H2O(l)} has $S^\circ = 70$ but \ce{H2O(g)} has 189
\si{\joule\per\kelvin\per\mole} --- the same substance, nearly triple, purely
because of the phase.}}

\problem[]{State the second law and explain why it refers to the universe
rather than the system.
\answerlines[4]{In any thermodynamically favored process the entropy of the
universe increases: $\Delta S_{\text{univ}} = \Delta S_{\text{sys}} +
\Delta S_{\text{surr}} > 0$. It must refer to the universe because the
system's own entropy is frequently allowed to fall --- water freezes,
ammonia forms, organisms build ordered molecules. Those are permitted
because the surroundings gain more entropy than the system loses.}}

\problem[]{$\Delta S_{\text{surr}} = -\Delta H/T$.}
\begin{parts}
  \item For an exothermic reaction, the sign of $\Delta S_{\text{surr}}$ is
        \answerblank[3cm]{positive}
  \item Why does the same quantity of heat raise the surroundings' entropy
        more at low temperature than at high?
        \answerlines[3]{The temperature sits in the denominator. Delivering
        a fixed amount of heat to cold surroundings is a large fractional
        increase in their energy dispersal; the same heat delivered to hot
        surroundings, which already hold a great deal of thermal energy,
        changes the dispersal comparatively little.}
  \item Use that to explain why water freezes below \SI{0}{\celsius} but
        melts above it, given that $\Delta H$ and $\Delta S$ for freezing
        barely change over that range.
        \answerlines[4]{Freezing has $\Delta S_{\text{sys}} < 0$ but is
        exothermic, so $\Delta S_{\text{surr}} = -\Delta H/T > 0$. As $T$
        falls that positive term grows, and below \SI{0}{\celsius} it
        outweighs the system's loss, making $\Delta S_{\text{univ}} > 0$.
        Above \SI{0}{\celsius} it no longer does, and melting becomes the
        favored direction instead.}
\end{parts}

\problem[]{For the combustion of carbon, \ce{C(gr) + O2(g) -> CO2(g)}:}
\begin{parts}
  \item Predict the sign of $\Delta S^\circ$ before calculating, and say
        why the prediction is hard here.
        \answerlines[2]{One mole of gas becomes one mole of gas, so
        $\Delta n_{\text{gas}} = 0$ and no large change is expected. With
        the dominant term absent, the sign cannot be predicted by
        inspection.}
  \item Calculate $\Delta S^\circ$. \workspace[2cm]
        \keyonly{\begin{apcanswer}$214 - (6 + 205) =
        \mathbf{+3}~\si{\joule\per\kelvin}$ --- very small, as
        predicted\end{apcanswer}}
  \item The reaction is strongly exothermic. What does that tell you about
        $\Delta S_{\text{univ}}$? \answerlines[2]{It must be strongly
        positive. The system contributes almost nothing, so essentially all
        of the increase comes from $\Delta S_{\text{surr}} = -\Delta H/T$,
        which is large because the reaction releases a great deal of heat.}
\end{parts}

\problem[]{A classmate says: ``Life creates order out of disorder, so living
things break the second law.'' Refute this.
\answerlines[4]{An organism is not an isolated system. It takes in
concentrated chemical energy and releases heat and simpler waste molecules
to its surroundings. The entropy gained by those surroundings exceeds the
entropy lost in assembling ordered biomolecules, so
$\Delta S_{\text{univ}}$ is still positive. The second law constrains the
universe, not any chosen part of it.}}

\problem[]{Rank by molar entropy, lowest first: \ce{H2O(s)}, \ce{H2O(l)},
\ce{H2O(g)}. Then explain what would change if the comparison were made at
\SI{500}{\kelvin} instead of \SI{298}{\kelvin}.
\answerlines[3]{\ce{H2O(s)} $<$ \ce{H2O(l)} $<$ \ce{H2O(g)} --- freedom of
motion and accessible volume both increase down the list. At
\SI{500}{\kelvin} every value would be larger, because raising the
temperature broadens the distribution of kinetic energies, but the ordering
would be unchanged.}}

\begin{spiralstrip}{Chapters 15--16 \textbullet\ spiral}
  \item pH $=$ p$K_a$ when the buffer ratio is:
        \answerblank[1.6cm]{1}
  \item For an \ce{MX2} salt, $K_{sp} = $ \answerblank[1.8cm]{$4s^3$}
  \item Common ion added $\Rightarrow$ solubility:
        \answerblank[1.8cm]{falls}
  \item Equivalence pH, weak acid $+$ strong base:
        \answerblank[2.2cm]{above 7}
  \item Does a common ion change $K_{sp}$? \answerblank[1.6cm]{no}
\end{spiralstrip}

\end{document}
